\documentclass{ieeeaccess}
\usepackage{cite}
\usepackage{amsmath,amssymb,amsfonts}
\usepackage{algorithmic}
\usepackage{graphicx}
\usepackage{textcomp}
\usepackage{booktabs}
\usepackage{multirow}
\usepackage{xcolor}
\usepackage{url}

\usepackage{bm}
\makeatletter
\AtBeginDocument{\DeclareMathVersion{bold}
\SetSymbolFont{operators}{bold}{T1}{times}{b}{n}
\SetSymbolFont{NewLetters}{bold}{T1}{times}{b}{it}
\SetMathAlphabet{\mathrm}{bold}{T1}{times}{b}{n}
\SetMathAlphabet{\mathit}{bold}{T1}{times}{b}{it}
\SetMathAlphabet{\mathbf}{bold}{T1}{times}{b}{n}
\SetMathAlphabet{\mathtt}{bold}{OT1}{pcr}{b}{n}
\SetSymbolFont{symbols}{bold}{OMS}{cmsy}{b}{n}
\renewcommand\boldmath{\@nomath\boldmath\mathversion{bold}}}
\makeatother

\def\BibTeX{{\rm B\kern-.05em{\sc i\kern-.025em b}\kern-.08em
    T\kern-.1667em\lower.7ex\hbox{E}\kern-.125emX}}

\newcommand{\TODO}[1]{\textcolor{red}{\textbf{[TODO: #1]}}}

%% ============================================================
\begin{document}
\history{Date of publication xxxx 00, 0000, date of current version xxxx 00, 0000.}
\doi{10.1109/ACCESS.2026.DOI}

\title{Temporal Pixel Masking for Reducing OCR Recovery from Short-Exposure Screen Photographs: A Single-UVC-Camera Feasibility Study}

\author{\uppercase{First A. Author}\authorrefmark{1}
\TODO{Author information to be filled}}

\address[1]{\TODO{Affiliation and address to be filled}}

\tfootnote{\TODO{Funding acknowledgment to be filled.}
This work includes a planned minimal-risk human-subject study. The authors'
institution does not maintain an ethics review board for behavioral studies
of this type; the study protocol follows the informed-consent and safety
procedures described in Section~\ref{sec:user_study}.}

\markboth
{Author \headeretal: Temporal Pixel Masking Against Short-Exposure Screen Photography}
{Author \headeretal: Temporal Pixel Masking Against Short-Exposure Screen Photography}

\corresp{Corresponding author: \TODO{Corresponding author email}.}

%% ============================================================
%% ABSTRACT
%% ============================================================
\begin{abstract}
Screen content can be silently captured by smartphones or smart glasses and subsequently processed by optical character recognition (OCR) or vision-language models to extract sensitive text. This paper proposes a temporal pixel mask that exploits differences in temporal integration: a cryptographically secure pseudorandom process decomposes each original frame into rapidly alternating complementary subframes, allowing human observers to integrate content across subframes while short-exposure cameras primarily record incomplete pixels. In a feasibility study conducted with a single USB Video Class (UVC) webcam across 9 distance/angle configurations and 10,575 captures, the best-of-engine OCR character recovery under the core short-exposure condition dropped from 94.1\% (unprotected) to 15.1\% for the readability-priority deployed profile and 5.0\% for the capture-hardened profile, with exact-match rates falling from 65.7\% to 0.4\% and 0\%, respectively. However, the capture-hardened profile still yields 47.9\% character recovery under video temporal averaging, and three commercial vision-language models can recover substantial portions of protected large-font short text. Accordingly, the validity claims in this paper are limited to conventional OCR short-exposure mitigation within the tested UVC link and are not extrapolated to other cameras, full-cycle integration, video reconstruction, or modern vision-language models. Color difference, structural similarity, and flicker metrics are numerical proxies; because the planned user study has not yet been conducted, actual readability and visual comfort remain unverified.
\end{abstract}

\begin{keywords}
Adversarial perturbation, anti-photography, camera capture, display privacy, OCR, persistence of vision, temporal masking, visual eavesdropping
\end{keywords}

\titlepgskip=-21pt
\maketitle

%% ============================================================
%% I. INTRODUCTION
%% ============================================================
\section{Introduction}
\label{sec:introduction}

\PARstart{S}{creens} are everywhere---in offices, banks, hospitals, and conference rooms. So too are capture devices: smartphones, smart glasses (e.g., Ray-Ban Meta)~\cite{b_wang2026}, and miniature cameras. An attacker needs only a single photograph of a screen to extract text via OCR or identify interface elements through object detection. This form of \emph{visual eavesdropping} is covert and easy to execute: it leaves no network intrusion traces and requires no physical contact with the target device. Controlled experiments and field studies confirm that such attacks are common in everyday settings, succeed at high rates, and frequently go unnoticed by victims~\cite{b_ponemon2016,b_eiband2017}. While these devices motivate this work, our physical evidence currently covers only a single UVC webcam in a screen-to-camera link; reproducibility on smartphones, smart glasses, and other camera ISP pipelines is left to future work.

\subsection{Existing Approaches and Their Limitations}

Current screen-privacy solutions each occupy a limited niche.
Physical privacy filters (e.g., commercial privacy screens) restrict viewing angles but cannot stop a camera positioned directly in front of the display. Software-based passive degradation (dimming, blurring, pixelation) sacrifices overall readability or depends on viewing geometry, and cannot specifically block on-axis capture---quantitative comparisons appear in the exploratory baseline experiment in \S\ref{sec:simulation}. Digital watermarking provides post-hoc accountability rather than pre-emptive blocking~\cite{b_fang2019}. Temporal re-encoding schemes designed for cinema anti-piracy (e.g., Kaleido~\cite{b_kaleido2015}) target degradation of continuous video recordings and do not address short-exposure single-frame capture of static sensitive content or machine recognition (\S\ref{sec:related}). These approaches respectively cover viewing-angle restrictions, post-hoc tracing, and continuous-recording disruption; on-axis short-exposure capture followed by machine recognition falls in the gap between them.

\textbf{The core tension}: any information degradation imposed on the camera tends to also reduce readability for the legitimate viewer. This paper asks whether display temporal structure can widen the gap between the two under short-exposure capture.

\subsection{Core Idea}

This paper exploits the difference between human temporal integration and camera exposure windows. We decompose each complete frame into $n$ randomly complementary subframes that alternate at high speed. When the exposure time is shorter than or close to one subframe duration, the camera primarily records an incomplete pixel subset, whereas the human observer integrates content across multiple subframes (Fig.~\ref{fig:concept}). Camera exposure is not instantaneous sampling. When long exposure or video reconstruction covers a full cycle, the complementary subframes can be re-summed; therefore, this paper makes no irrecoverability guarantee for arbitrary exposure conditions.

Short exposure is modeled as a relevant first attacker tier for opportunistic capture: an observer can attempt one brief snapshot without first calibrating exposure or keeping the camera aimed at the screen for multiple frames. This is not a claim that every default smartphone selects a subframe-length shutter---the S600 experiments use fixed, calibrated exposures rather than smartphone auto-exposure measurements. A capable attacker can deliberately integrate longer or record video, so those modes form explicit higher attacker tiers and a security boundary rather than an excluded case.

\begin{figure}[!t]
\centering
\includegraphics[width=\columnwidth,viewport=0 0 252 168,clip]{figures/visio/figure1_concept_threat/final/figure1_concept_threat.pdf}
\caption{Threat scenario and operating boundary of temporal pixel masking. (a) A nearby camera can capture sensitive on-screen content without accessing the host system. (b) Human vision integrates complementary subframes, while a short camera exposure can intersect only part of a display cycle and leave OCR with an incomplete fragment. The camera path is conceptual: it does not assume zero-duration sampling or a fixed human integration time, and long exposure or multi-frame reconstruction can recover substantially more information.}
\label{fig:concept}
\end{figure}

\subsection{Contributions}

The main contributions of this paper are:
\begin{enumerate}
\item \textbf{A temporal pixel mask implementation targeting conventional OCR under short-exposure capture}: Building on prior temporal masking ideas, we implement CSPRNG-driven pixel-level complementary subframes, optional complementary adversarial noise, and a GPU rendering prototype, with an explicit acknowledgment that effectiveness depends on camera exposure windows and the attacker's recognition pipeline.

\item \textbf{Multi-geometry real-capture OCR feasibility evaluation on a single UVC camera with layered evidence}: Using one eMeet SmartCam S600 USB webcam, we aggregate 10,575 captures across 3 distances and 3 angles. This corpus comprises component ablations, parameter sweeps, and attack evaluations rather than a balanced full-factorial design; the total count conveys within-link coverage, and the primary validity conclusions draw only from the core short-exposure OCR subset: 15.1\% character recovery (0.4\% exact) for the deployed profile and 5.0\% (0\% exact) for the capture-hardened profile.

\item \textbf{Cross-task stress tests and VLM attack boundary probes}: Beyond the OCR main experiment, we report results for 4 object detectors on COCO val2017 (150 real-capture images, plus 8 for simulation pipeline diagnostics), MOT17 detection/association (5,316 simulated frames and 450 real-capture still frames), and 3 commercial VLMs at the 0$^{\circ}$ on-axis position across three distances. The VLM experiment serves as negative boundary evidence: conventional OCR protection does not automatically transfer to modern end-to-end vision recognition. Conclusions are bounded by the evaluated models, positions, and tasks.

\item \textbf{An honest security-frontier analysis}: We quantify residual leakage under integration attacks: ideal aligned full-cycle integration recovers 95.2\% in digital simulation; real-capture video temporal averaging against the readability-priority deployed profile still recovers 71.1\% char. The capture-hardened profile reduces real-capture video leakage to 47.9\% char / 5.6\% exact, with an explicit statement that this attack class remains the fundamental boundary of the scheme. Real-capture VLM probes further demonstrate that short-exposure protection against large-font short snippets fails for modern VLMs: at 0.5\,m Qwen3-VL reaches 77.8\% exact, and at 1.5\,m the three models retain 33.3--47.2\% short-exposure exact. This result narrows the contribution from ``defeating all screen-reading attacks'' to ``conventional OCR short-exposure mitigation and boundary characterization of its failure against VLMs.''

\item \textbf{Security--usability evaluation framework}: We report FPI, $\Delta E$, SSIM, and other proxy metrics together with Pareto fronts, explicitly distinguish digital integration proxies from actual user experience, and specify a controlled within-subject protocol for subsequent usability validation.
\end{enumerate}

The remainder of this paper is organized as follows. Section~\ref{sec:related} reviews prior display-privacy and temporal-modulation work. Section~\ref{sec:threat} defines attacker tiers and protection goals. Section~\ref{sec:method} presents the masking design, and Section~\ref{sec:evaluation} reports OCR, VLM, cross-task, usability-protocol, and security-frontier evaluations. Sections~\ref{sec:discussion} and~\ref{sec:conclusion} discuss the deployment boundary and conclude the paper.


%% ============================================================
%% II. RELATED WORK
%% ============================================================
\section{Related Work}
\label{sec:related}

\subsection{Screen Visual Eavesdropping Threats}

Remote screen reading attacks (telescopes, telephoto lenses)~\cite{b_backes2009}, deep-learning reconstruction of HDMI electromagnetic leakage~\cite{b_fernandez2024}, covert capture capabilities of smart glasses~\cite{b_wang2026}, and indirect eavesdropping through reflections and refractions~\cite{b_raguram2011} have all been documented. A controlled experiment by the Ponemon Institute found that 91\% of visual hacking attempts succeed within 15~minutes~\cite{b_ponemon2016}. Eiband et al.~\cite{b_eiband2017} surveyed 174 participants to reveal the prevalence of shoulder surfing in everyday settings and the coping strategies users adopt. Tang and Shin~\cite{b_tang2023} proposed Eye-Shield, which uses the front-facing camera to detect bystanders and blur sensitive screen regions in real time, though it relies on user-side hardware and only provides localized masking. This paper focuses on display-side mitigation that does not require attacker-side cooperation, without claiming that information acquisition can be prevented under all exposure and reconstruction conditions.

\subsection{Display Privacy Protection}

Physical privacy filters restrict viewing angles to protect against lateral observers but cannot block an on-axis camera. Chen et al. proposed HideScreen, a display-side shoulder-surfing defense that uses grid-based rendering for sensitive regions and suppresses low-frequency visual cues so screen content blends with the background beyond a designed viewing distance~\cite{b_hidescreen2019}. Its target is human shoulder surfing in the spatial domain; it does not address short-exposure camera capture or machine recognition in the temporal sampling domain. Zhu et al.~\cite{b_zhu2017} proposed LiShield, which uses smart LED modulated illumination to interfere with rolling-shutter cameras for physical-scene protection, but the scheme acts on ambient lighting rather than displayed content. Visual cryptography splits a secret image into shares that combine upon overlay~\cite{b_naor1994}, inspiring temporal decomposition ideas, but its original scheme targets printed media rather than high-refresh displays. Digital watermarking~\cite{b_zhong2023} can embed tracking identifiers---recent advances include screen-shooting-resilient watermarks~\cite{b_fang2019}---but watermarking inherently provides post-hoc attribution rather than pre-emptive blocking. Software-based DRM prevents digital screenshots but not physical photography.

Screen-camera visible light communication (VLC)~\cite{b_nguyen2016,b_tran2021,b_ghiasi2024} similarly exploits human temporal integration to embed information imperceptible to the eye yet decodable by cameras. Our method is VLC's mirror image: VLC aims to let cameras \emph{read} hidden information, whereas we aim to prevent cameras from \emph{reading} displayed content.

Yerazunis and Carbone~\cite{b_yerazunis2001} first proposed using temporal image masking to enhance display privacy. Wu and Zhai~\cite{b_wu2013} subsequently systematized temporal psychovisual modulation, and Gao et al.~\cite{b_gao2014} further applied it to information-security displays. Zhang et al.~\cite{b_kaleido2015} proposed Kaleido, which re-encodes video frames on standard-refresh-rate screens using multi-frame decomposition to exploit the channel difference between display-eye and display-camera paths, making films ``viewable but not recordable''; Yang et al.~\cite{b_rainbow2018} extended this anti-piracy direction with Rainbow, covering mobile-camera filming in physical environments. These schemes aim to degrade continuous recording quality to counter video piracy, evaluating via PSNR, SSIM, and subjective quality scores on recorded footage---metrics that are non-interchangeable with the OCR, object detection, and VLM machine recovery rates used here, and none address short-exposure single-frame capture of static sensitive content coupled with machine recognition. Mechanistically, Kaleido's decomposition operates in continuous color space with chrominance-complementary frames and luminance variation to degrade recorded viewing experience, relying on perceptually complete mixing within a cycle for the human viewer; our method employs CSPRNG-driven discrete pixel-level mutually exclusive partitioning, and the capture-hardened profile further deliberately abandons strict completeness to resist integration attacks (\S\ref{sec:hardened}). Lim~\cite{b_lim2007} used persistence of vision to generate anti-screenshot keyboards in a browser, where a single digital screenshot captures only partial patterns. Thus, temporal partitioning per se is not novel to this paper. Our incremental contributions are: (1)~pixel-level CSPRNG complementary partitioning with a hardened profile targeting physical cameras and conventional OCR under short exposure; (2)~10,575 real captures covering 9~geometric conditions on a single UVC webcam, with the core short-exposure OCR subset serving as feasibility evidence; (3)~object detection, object association, and VLM probes as cross-task stress tests rather than claims of general visual protection; and (4)~quantitative boundary analysis of full-cycle averaging and VLM failure.

\subsection{Adversarial Perturbations and Physical-World Robustness}

Digital adversarial perturbations~\cite{b_goodfellow2015,b_madry2018,b_carlini2017} and physical-world adversarial patches~\cite{b_eykholt2018,b_thys2019,b_xu2020} have been extensively surveyed~\cite{b_wei2024}. Athalye et al.~\cite{b_athalye2018} achieved the first cross-viewpoint-robust 3D physical adversarial objects via Expectation over Transformation (EoT), though their approach still requires white-box access to the target model. A significant digital-to-physical robustness gap persists: pixel-level digital perturbations are substantially weakened after camera sampling and JPEG compression~\cite{b_eykholt2018}. Physical adversarial patches exhibit cross-viewpoint robustness but require perturbation patterns pre-generated for specific classifiers~\cite{b_thys2019,b_xu2020}, with limited transferability to unseen models~\cite{b_zhao2023}. The screen-camera link also introduces moir\'{e} artifacts, which Niu et al.~\cite{b_niu2021} demonstrated can themselves act as adversarial perturbations. Our method does not add perturbations to content; instead, it performs temporal partitioning at the display layer. Physical photography is inherently a sampling process, and the defense is embedded in the display itself, requiring no assumptions about the attacker's recognition model.


%% ============================================================
%% III. THREAT MODEL
%% ============================================================
\section{Threat Model}
\label{sec:threat}

\subsection{Protected Assets}

Sensitive text and visual information displayed on screen: passwords, financial data, personally identifiable information (PII), source code, confidential documents, and similar content.

\subsection{Attacker Capabilities}

Attack capabilities do not form a strict monotonic hierarchy but are jointly determined by exposure, geometry, temporal sampling, and post-processing capabilities. Results are reported across the following scope:
\begin{itemize}
\item \textbf{Primary protection scope}: Opportunistic short-exposure snapshots and per-frame recognition by conventional OCR attackers. Camera exposure time is shorter than or close to one subframe duration; the attacker may choose among OCR engines but does not possess observations covering a full display cycle. This tier represents a brief single-shot attempt without exposure calibration or sustained aiming, not the strongest rational attacker.
\item \textbf{Evaluation boundary}: Long exposure, video temporal averaging after continuous recording, off-axis capture, phase search and channel reconstruction given a known period, object detection/association tasks, and end-to-end recognition by commercial vision-language models (\S\ref{sec:vlm}). This paper quantifies these attack or task-transfer risks but interprets capture-hardened results only as mitigation, not as a complete blocking guarantee.
\item \textbf{Not yet adequately covered}: Additional camera/display combinations, precise rolling-shutter row alignment, real-world temporal super-resolution, learning-based reconstructors trained on display sequences, and long-term adaptive attacks.
\end{itemize}

\subsection{Attacker Knowledge}

The attacker \emph{knows} the temporal masking algorithm, the subframe count~$n$, and candidate period parameters, but does \emph{not know} the current cycle's mask seed or phase starting point. The CSPRNG prevents fixed spatial patterns and cross-cycle repetition but does not provide a cryptographic guarantee against full-cycle averaging attacks. Once an attacker obtains and registers a complete cycle, an unknown seed cannot prevent linear reconstruction.

\subsection{Protection Goals}

This paper defines tiered design goals under short-exposure conditions by protection profile: the capture-hardened profile targets conventional OCR character recovery $<5\%$ with exact-match $=0\%$ as its strict pass criterion; the readability-priority deployed profile targets character recovery reduction $\geq$80\% relative to the unprotected baseline with exact-match $<1\%$. Changes in object detection mAP50 are used to explore task-transfer risk beyond text and are not primary pass criteria. These thresholds guide parameter selection and profile trade-offs but do not constitute formal security guarantees against all attack modes. The thresholds are defined for conventional OCR attackers; VLM attacker results are not included in pass/fail determination but are reported separately as security boundaries (\S\ref{sec:vlm}). On the usability side, FPI, $\Delta E$, and SSIM serve as modeled proxies. The planned user study has not yet been conducted, so this paper does not claim ``imperceptible to the human eye'' based on these metrics.


%% ============================================================
%% IV. SYSTEM DESIGN
%% ============================================================
\section{System Design}
\label{sec:method}

\subsection{Temporal Subframe Decomposition}

Given a normalized original frame $\mathbf{I}\in[0,1]^{H\times W\times C}$, we generate binary masks $\mathbf{M}_1,\ldots,\mathbf{M}_n$ satisfying $\sum_{k=1}^{n}\mathbf{M}_k=\mathbf{1}$, and set $\mathbf{I}_k=\mathbf{I}\odot\mathbf{M}_k$. This yields:
\begin{itemize}
\item \textbf{Digital-domain completeness}: $\sum_{k=1}^{n}\mathbf{I}_k = \mathbf{I}$;
\item \textbf{Mutual exclusivity}: each pixel is illuminated in exactly one of the $n$ time slots.
\end{itemize}
A distinction between ``summation'' and ``temporal averaging'' is necessary. On an ideal linear display with equal slot durations and equal peak luminance, the average radiance over one cycle is $\mathbf{I}/n$, not $\mathbf{I}$. Digital-domain completeness means subframes can be re-summed; it does not imply that real display luminance has been losslessly preserved.

Masks are generated using a ChaCha20~\cite{b_chacha2008} CSPRNG and the Fisher--Yates shuffle~\cite{b_knuth1997}, randomly assigning each pixel to one subframe. The implementation applies a chi-square uniformity check on generated slot counts, re-sampling the seed upon failure. Random seeds reduce fixed patterns and cross-cycle correlation but do not change the fact that a full cycle can be linearly summed. Fig.~\ref{fig:pipeline} illustrates the complete pipeline from original frame to protected display.

\begin{figure*}[!t]
\centering
\includegraphics[width=\textwidth,viewport=0 0 515.52 206.16,clip]{figures/visio/figure2_method_pipeline/final/figure2_method_pipeline.pdf}
\caption{Temporal pixel masking system pipeline from source content to protected display. The input frame is decomposed into CSPRNG-assigned complementary pixel masks; optional complementary OCR-oriented noise, profile-dependent stripe/glyph overlays, and an optional partial inversion frame are then inserted before high-refresh playback. A human observer is expected to integrate multiple slots into a readable view, whereas a short-exposure camera may capture only a sparse subframe. Long exposure and video temporal averaging can accumulate multiple slots and are therefore evaluated separately as boundary attacks rather than assumed away.}
\label{fig:pipeline}
\end{figure*}

\subsection{Complementary Adversarial Noise}

We overlay adversarial noise $\mathbf{N}_1, \ldots, \mathbf{N}_n$ on the subframes, satisfying:
\begin{equation}
\sum_{k=1}^{n}\mathbf{N}_k = \mathbf{0}
\label{eq:noise_complement}
\end{equation}
Equation~(\ref{eq:noise_complement}) holds strictly only in the ideal linear domain prior to summation, clipping, and the display transfer function. Prior work has shown that small text-image perturbations can significantly mislead OCR systems~\cite{b_song2018}, providing a basis for OCR-oriented noise design. Noise directions are generated by FGSM-style~\cite{b_goodfellow2015} proxy gradients with an amplitude bound of $\varepsilon=8/255$, then decomposed complementarily in the temporal domain. Because displays cannot emit negative radiance, the implementation uses a pedestal level and numerical clipping. Consequently, after gamma correction, quantization, clipping, and camera imaging, the noise is not guaranteed to cancel physically.

\textbf{Component role}: The random dot-pattern mask is the primary mechanism. In controlled digital experiments, noise can suppress residual structure under weak masking or repair attacks. However, in real-capture short-exposure results, mask+noise (25.6\%) is higher than mask-only (19.2\%), and the Strong profile with weak stripe/glyph overlay is also slightly higher than mask-only (20.7\% versus 19.2\%) under the best-of-engine caliber. Therefore, this paper does not present adversarial noise or weak-amplitude overlay as universally monotonic enhancements across all physical conditions, but rather as optional/profile-dependent components whose benefit depends on the attack, recognizer selection, and imaging chain.

\subsection{Luminance Compensation}

Each pixel is illuminated only once per $n$-slot cycle. In theory, if the display could boost light output by a factor of $n$ in the linear-light domain during the active slot, the cycle-averaged luminance would approach the original. However, the current PoC does not implement per-slot dynamic backlight coordination, nor does it use a photometer to verify absolute luminance. The luminance alignment, CIEDE2000 $\Delta E_{00}$~\cite{b_ciede2000}, and SSIM results in this paper all come from digital integration/normalization models and are used only for cross-configuration comparison, not as evidence of real-panel optical equivalence or user comfort.

\subsection{Timing and Bandwidth}

For a basic cycle containing only $n$ complementary subframes, taking 60\,Hz as the minimum cycle rate requires $f_r\geq60n$; at $n=4$ this corresponds to 240\,Hz. If the cycle additionally includes one inversion frame, the requirement becomes $f_r\geq60(n+1)$, and at 240\,Hz the actual cycle rate drops to only 48\,Hz. Thus, 240\,Hz meets the basic four-subframe configuration but not the inversion-frame configuration's 60\,Hz cycle target. FPI is a modeled flicker proxy based on cycle rate and spatial pooling; lower values indicate lower predicted flicker risk rather than a clinical or user-study threshold. The spatial phase variation of random masks may reduce local synchronized modulation, but human flicker perception is additionally influenced by luminance, contrast, retinal eccentricity, and eye movements~\cite{b_cai2024,b_davis2015}.

\subsection{Capture-Hardened Profile}
\label{sec:hardened}

The basic linear temporal mask leaks under full-cycle temporal averaging attacks (see \S\ref{sec:security_frontier}). To address this frontier, nonlinear enhanced profiles are introduced. Unless otherwise stated, the real-capture profiles in this subsection use $n=4$ base temporal masking with complementary adversarial noise enabled. Thus the ``Mask + noise'' row in Table~\ref{tab:real_ocr_main} is the base mask+noise condition without stripe/glyph anti-OCR overlay; the Strong row adds the weak stripe/glyph overlay to that base, the Deployed row adds a weak inversion frame to Strong, and the capture-hardened row uses stronger stripe/glyph settings plus the same weak inversion frame.
\begin{itemize}
\item \textbf{Strong (anti-OCR) profile}: base mask+noise with mask cell size of 1~pixel, stripe width of 10~pixels, stripe amplitude $\alpha_s=0.10$, and glyph amplitude $\alpha_g=0.12$, without an inversion frame.

\item \textbf{Deployed profile}: the Strong profile plus an $\alpha=0.2$ inversion frame, selected as the readability-priority configuration.

\item \textbf{Capture-hardened profile}: base mask+noise with mask cell size increased to 2~pixels, stripe width of 6~pixels, stripe and glyph amplitudes of 0.42 and 0.55 respectively, plus an $\alpha=0.2$ inversion frame.
\end{itemize}

The nonlinear profiles depart from strict completeness ($\sum\mathbf{I}_k \neq \mathbf{I}$), producing visible stripes and grain. This trade-off is the opposite of Kaleido-class temporal re-encoding: the latter relies on perceptually complete mixing within a cycle to preserve viewing experience, meaning full-cycle integration inherently reconstructs an approximation of the original---consistent with its anti-piracy goal. In the static-content capture threat addressed here, completeness is precisely the entry point for integration attacks; the hardened profile therefore makes breaking strict completeness an active design trade-off, accepting visible stripes and grain in exchange for robustness to full-cycle integration. Its security benefit and unmeasured readability cost must be reported separately.

\subsection{Partial Inversion Frame}

A long-exposure attack occurs when the exposure window covers a full display cycle ($n=4$ at 240\,Hz yields a basic cycle of approximately 16.7\,ms; with an inversion frame the $n+1$ slot cycle is approximately 20.8\,ms; \S\ref{sec:setup} reports measured long-exposure times of 31.25 and 125\,ms). Such attacks approximate the original image through sensor integration. A direct countermeasure inserts a \textbf{partial inversion frame} $\alpha \cdot (255 - \mathbf{I})$ after each $n$-subframe cycle, introducing a luminance component complementary to the original into the long-exposure integral and reducing net signal contrast. Higher $\alpha$ may reduce long-exposure recovery, but also alters the luminance and contrast of the digital integrated view.

A real-capture inversion strength ablation (Table~\ref{tab:inversion_ablation}; capture and evaluation pipeline per \S\ref{sec:real_capture}) shows that weak inversion does not improve long-exposure character recovery: $\alpha=0$ yields 68.6\% and $\alpha=0.2$ yields 67.0\% (best-of-engine), with 95\% confidence intervals overlapping substantially across $\alpha\in[0,0.5]$. The ablation holds the Strong overlay fixed ($\alpha_s=0.10$, $\alpha_g=0.12$) and varies only inversion strength, but it uses the five-item account/code/CET6 subset rather than the 12-item main-profile corpus. Its $\alpha=0.2$ long-exposure value ($N=135$) is therefore not a repeated estimate of the deployed row's 60.9\% ($N=324$). Only near-full inversion ($\alpha=1.0$) reduces recovery to 18.1\%, but at that setting the digital integrated view is severely degraded. Under video temporal averaging, weak inversion similarly maintains approximately 64--73\% character recovery. Therefore, weak inversion frames cannot independently resist integration attacks. The deployed profile's choice of $\alpha=0.2$ represents a modeled readability-priority design decision, not a user-study-validated optimum.

\begin{table}[!t]
\caption{Inversion Frame Strength $\alpha$ Ablation on the Fixed Strong Overlay ($\alpha_s=0.10$, $\alpha_g=0.12$): Real-Capture Best-of-Engine Character Recovery (Mean and 95\% Bootstrap CI, Five Account/Code/CET6 Items Pooled Across 9 Geometric Conditions; $N=135$ per Setting for Long Exposure, $N=45$ for Video Temporal Averaging). This Five-Item Subset Differs from the 12-Item Deployed Rows in Table~\ref{tab:real_ocr_main}.}
\label{tab:inversion_ablation}
\centering
\small
\begin{tabular}{lcc}
\toprule
$\alpha$ & Long-exp.\ char (\%) & Video avg.\ char (\%) \\
\midrule
0.0 & 68.6~[63.1, 73.9] & 72.7~[63.3, 81.1] \\
0.2 & 67.0~[61.3, 72.7] & 69.8~[59.9, 78.7] \\
0.3 & 72.0~[66.2, 76.9] & 64.3~[54.0, 73.6] \\
0.5 & 61.7~[55.6, 67.7] & 66.4~[56.5, 75.2] \\
1.0 & 18.1~[14.0, 22.0] & 28.4~[20.5, 36.9] \\
\bottomrule
\end{tabular}
\end{table}

Each display cycle thus contains $n+1$ time slots. At $n=4$ and 240\,Hz the cycle rate is 48\,Hz, below this paper's 60\,Hz design target and potentially producing perceptible flicker; no fixed ``flicker-free threshold'' applies universally across viewing conditions.

\subsection{Implementation Scope}

The current PoC operates at the application/algorithm layer, using Python and GPU rendering to generate display sequences. It does not include OS driver-level injection (DXGI / KMS-DRM / CoreDisplay) or per-slot dynamic backlight control. Accordingly, this paper validates the display sequence and its captured results, without claiming production-level system integration for the current prototype.


%% ============================================================
%% V. EXPERIMENTAL EVALUATION
%% ============================================================
\section{Experimental Evaluation}
\label{sec:evaluation}

\subsection{Experimental Setup}
\label{sec:setup}

\textbf{Hardware platform}: Experiments ran on a desktop workstation with an Intel Core i7-8700K (3.70\,GHz), 32\,GB RAM, and an NVIDIA TITAN Xp (12\,GB VRAM). The display canvas resolution was $1920 \times 1080$ at a nominal refresh rate of 240\,Hz. This refresh rate satisfies only the basic $n=4$ subframe cycle at $f_r/n=60$\,Hz. With an additional inversion frame per cycle, the actual cycle frequency is $240/(n+1)=48$\,Hz. Thus, the visual comfort of the inversion configuration cannot be inferred from ``240\,Hz high refresh rate'' alone and still requires user-study validation.

\textbf{Capture device}: The physical capture experiments are positioned as a single-UVC-camera feasibility study using an eMeet SmartCam S600 USB webcam (UVC standard driver). The playback canvas was $1920 \times 1080$; camera output underwent perspective correction and content cropping before being fed to recognizers. Merged experimental metadata indicate short-exposure and video-frame exposure times of 0.49 or 3.91\,ms, long-exposure times of 31.25 or 125\,ms, and video capture at 60\,fps. The camera was placed at 0.5\,m, 1\,m, and 1.5\,m from the display, and at each distance rotated to $0^{\circ}$, $15^{\circ}$, and $30^{\circ}$, yielding 9~geometric configurations. This setup validates feasibility within the current UVC screen-camera link; ambient illuminance was not recorded, and the study does not cover smartphones, smart glasses, different sensors, different lenses, or other camera ISPs. Results are therefore not extrapolated to these devices or lighting conditions.

\textbf{Test content}: At each of the 9 geometric positions, the real-capture OCR experiment used a fixed set of 12 content items: 11 selected synthetic items and one CET6 (College English Test Band~6) document (portrait $900 \times 1280$). Three captures per item give 36 captures per profile--attack condition at each position and $12\times3\times9=324$ across all positions before the disclosed deployed-profile repeats. By contrast, the software-simulated OCR experiment used the complete 120-sample synthetic corpus (12 content templates $\times$ 10~variants, mapped to 9 category labels covering passwords, credentials, code, tables, paragraphs, URLs, etc.). Real-capture detection used a 150-image COCO val2017~\cite{b_coco2014} subset; tracking used 450~frames from the MOT17-09-FRCNN~\cite{b_mot16} sequence. Simulated detection used only 8~images for pipeline diagnostics; simulated tracking covered 5,316~frames. Accordingly, real-capture short-exposure OCR is the primary feasibility evidence; detection, tracking, and simulated detection results serve to explore task-transfer risk beyond text, and the 8-image simulated detection is not claimed as a full-scale evaluation.

\textbf{Evaluation metrics}: OCR character recovery is defined as $\max(0,1-\mathrm{CER})$. Exact-match is determined after whitespace removal and case-insensitive comparison of the full text. Sensitive-token recovery counts the proportion of reference-text digits, account numbers, URL/path strings, and key-like fields that are recovered in full. Leak rate is the proportion of samples with character recovery $\geq$20\%. Detection uses mAP, mAP50, and AR; tracking uses HOTA and IDF1, with MOTA/MOTP additionally reported in simulation. Section~\ref{sec:user_study} defines the separate usability protocol and metrics.

For the real-capture OCR main results, we first take the maximum recovery across three engines per capture to model an attacker choosing the best recognizer, then estimate 95\% bootstrap confidence intervals with 2,000 percentile resamples using captures as the resampling unit and a fixed random seed of 20260612. Confidence intervals for core conditions are reported inline in \S\ref{sec:real_capture}; ablation intervals appear in Table~\ref{tab:inversion_ablation}. Because repeated measurements exist for the same content and geometric conditions, these intervals describe empirical uncertainty within the current corpus and capture setup, not independent population inference. Detection and tracking results have not undergone significance testing; differences are reported descriptively.

Digital usability proxies: $\Delta E_{00}$ is the mean CIEDE2000 color difference between the original image and the full-cycle digitally reconstructed image; SSIM is their mean structural similarity. The normalized mutual information proxy is defined as $I(X;Y)/H(X)$, where $X$ is the original grayscale and $Y$ is a single-subframe grayscale, estimated with 64 histogram bins; lower values indicate less information retained by a single subframe. FPI uses a simplified model in our implementation: let single-pixel modulation depth $d=1-1/n$, spatial pooling pixel count $N=625$, and pixel illumination frequency $f_p=f_r/n$. When $f_p\geq60$\,Hz, $\mathrm{FPI}=dN^{-1/2}(60/f_p)^2$; below 60\,Hz the formula uses $d(1-f_p/60)$. Note that this piecewise form is discontinuous at $f_p=60$\,Hz: in a narrow band just below 60\,Hz the low-frequency branch yields values smaller than the high-frequency branch at 60\,Hz, contradicting the intuition that lower frequency means worse flicker. The 12~configurations swept in this paper (with $f_p$ grid points from 18--180\,Hz) do not land in this critical narrow band, so configuration ranking is unaffected. This formula is not an IEEE~1789 standard metric or clinical threshold; it serves only for within-paper configuration ranking.

% ----------------------------------------------------------
\subsection{Real-Capture OCR Experiment}
\label{sec:real_capture}

The real-capture OCR corpus comprises 10,575 images, with 1,175 per geometric position, all from the same eMeet SmartCam S600 UVC webcam. The corpus mixes component ablations, parameter sweeps, and main attack conditions, and is therefore not a balanced ``9~geometries $\times$ 6~main profiles $\times$ 3~attacks'' full-factorial design; stripe, glyph, and inversion settings are separate parameter sweeps, not a seventh main profile. One asymmetry in sample counts requires disclosure: the deployed profile performed two rounds of same-parameter captures for 5 of the 12 content items at each position (2 account credentials, 2 code snippets, 1 CET6 document); all other profiles used one round. Its short-exposure/video sample counts are therefore 459/153 rather than 324/108; both rounds are included in analysis. This paper uses the total count to convey within-link physical capture coverage while using the core attack-condition subset as primary feasibility evidence; these results do not constitute cross-camera, cross-smartphone, or cross-smart-glasses statistical validation. Three OCR engines---Surya~\cite{b_surya2024}, EasyOCR~\cite{b_easyocr2020}, and Tesseract~\cite{b_smith2007}---were each applied, producing 31,725 engine-level records. Recent vision-language models have demonstrated strong OCR capabilities~\cite{b_shi2023}; the real-capture VLM probes in \S\ref{sec:vlm} confirm that protection effective against conventional OCR cannot be extrapolated to VLM attackers. Table~\ref{tab:real_ocr_main} reports the core condition subset with per-capture best-of-engine. Table~\ref{tab:real_ocr_engine} provides per-engine results under short exposure. Fig.~\ref{fig:montage} shows qualitative comparisons of digital integrated views and short/long-exposure real captures; Fig.~\ref{fig:bar_attack} summarizes recovery rates by profile $\times$ attack mode.

\begin{table*}[!t]
\caption{Single-UVC Real-Capture OCR Core Conditions (Best-of-Engine; $N$ = Independent Captures per Row; Pooled Across 9 Geometric Conditions)}
\label{tab:real_ocr_main}
\centering
\begin{tabular}{lllcccc}
\toprule
Profile & Attack mode & $N$ & Char recovery & Exact match & Sensitive token & Leak rate \\
\midrule
\multirow{3}{*}{Original (unprotected)}
 & short & 324 & 94.1\% & 65.7\% & 98.4\% & 99.4\% \\
 & video:temporal\_mean & 108 & 79.9\% & 61.1\% & 85.6\% & 86.1\% \\
 & long & 324 & 47.3\% & 18.8\% & 95.7\% & 58.6\% \\
\midrule
\multirow{3}{*}{Mask only}
 & short & 324 & 19.2\% & 0.3\% & 30.8\% & 30.6\% \\
 & video:temporal\_mean & 108 & 79.0\% & 48.1\% & 84.0\% & 88.0\% \\
 & long & 324 & 67.2\% & 35.5\% & 94.0\% & 75.0\% \\
\midrule
\multirow{3}{*}{Mask + noise}
 & short & 324 & 25.6\% & 2.8\% & 34.4\% & 36.1\% \\
 & video:temporal\_mean & 108 & 79.2\% & 49.1\% & 83.9\% & 88.9\% \\
 & long & 324 & 68.0\% & 39.8\% & 95.7\% & 76.2\% \\
\midrule
\multirow{3}{*}{Strong (anti-OCR)}
 & short & 324 & 20.7\% & 0.6\% & 33.7\% & 32.4\% \\
 & video:temporal\_mean & 108 & 78.6\% & 50.9\% & 83.3\% & 88.9\% \\
 & long & 324 & 61.6\% & 39.8\% & 95.7\% & 71.3\% \\
\midrule
\multirow{3}{*}{Deployed}
 & short & 459 & 15.1\% & 0.4\% & 24.0\% & 20.5\% \\
 & video:temporal\_mean & 153 & 71.1\% & 42.5\% & 78.5\% & 85.0\% \\
 & long & 324 & 60.9\% & 36.4\% & 96.5\% & 69.1\% \\
\midrule
\multirow{3}{*}{Capture-hardened}
 & short & 324 & 5.0\% & 0.0\% & 6.5\% & 3.7\% \\
 & video:temporal\_mean & 108 & 47.9\% & 5.6\% & 45.6\% & 76.9\% \\
 & long & 324 & 9.3\% & 0.0\% & 34.9\% & 14.2\% \\
\bottomrule
\end{tabular}
\end{table*}

\textbf{Key findings}:

\begin{enumerate}
\item \textbf{Substantial recovery reduction under short exposure}: Unprotected screen photographs yield 94.1\% OCR character recovery (95\% CI [93.1, 95.0]; exact 65.7\%). The strong (anti-OCR) profile without inversion frame achieves 20.7\% char / 0.6\% exact; adding an $\alpha=0.2$ weak inversion frame (deployed profile) reduces this to 15.1\% (CI [13.3, 17.1]; exact 0.4\%), an 83.9\% relative reduction from the unprotected condition and meeting the deployed-profile goal defined in \S\ref{sec:threat}. The descriptive difference between the two profiles suggests the weak inversion frame provides additional suppression under the current short-exposure setting, but the deployed profile includes additional repeated captures and has $N=459$ rather than 324, so the 5.6-percentage-point gap is not interpreted as a strict paired causal effect. The capture-hardened profile further reduces to 5.0\% (CI [4.3, 5.8]; exact 0\%); its exact-match satisfies the strict pass criterion, while character recovery sits at the 5\% boundary with the confidence interval straddling the threshold---this paper does not claim it fully satisfies the strict criterion. These results support short-exposure mitigation feasibility within the current eMeet S600 UVC camera, exposure, and geometry range, but do not constitute guarantees for arbitrary camera settings, smartphone cameras, or smart glasses.

\item \textbf{Integration attacks are the residual frontier; the readability-priority profile is vulnerable here}: Both long exposure and video temporal averaging exploit full-cycle linear complementarity. The readability-priority deployed profile still recovers 60.9\% char / 36.4\% exact under long exposure (sensitive-token recall 96.5\%) and 71.1\% char / 42.5\% exact under video temporal averaging. Notably, the unprotected long-exposure baseline is only 47.3\% char, lower than the deployed profile's 60.9\%, indicating that fixed long-exposure/camera parameters produce nonmonotonic effects across display conditions. This set of absolute values cannot be interpreted as a universal protection benefit. Weak inversion ($\alpha=0.2$) is likewise insufficient to independently resist temporal integration. The weak Strong overlay also should not be read as a measured integration-attack defense: mask-only and Strong are nearly identical under video temporal averaging (79.0\% versus 78.6\%) and differ only modestly under long exposure (67.2\% versus 61.6\%), with overlapping bootstrap intervals in the archived summaries. The monotone synthetic trend in Fig.~\ref{fig:tradeoff} therefore does not imply a measurable real-capture integration gain at weak overlay amplitude; the real-capture integration reduction appears only at capture-hardened amplitudes, and even there leakage remains substantial. Accordingly, the primary validity conclusions are confined to the evaluated single-frame short-exposure threat model; integration attacks are reported separately as a security boundary.

\item \textbf{The capture-hardened profile reduces recovery under current integration attacks but retains leakage and an unresolved perception cost}: Under the current setup, this profile yields 9.3\% char (CI [7.9, 10.9]) / 0\% exact for long exposure and 47.9\% char (CI [41.8, 54.0]) / 5.6\% exact (CI [1.9, 10.2]) for video temporal averaging. The video condition still shows substantial leakage and should not be read as an attack closure. Additionally, 47.9\% is the pooled value across 9~positions with large inter-position variation: 0.5\,m at three angles yields 50.7\%/0\%/33.8\%; 1.0\,m yields 61.1--69.4\%; 1.5\,m yields 48.4--55.5\%. The 0\% position (0.5\,m/$15^{\circ}$) is the only one calibrated with the shorter 0.49\,ms video-frame exposure (\S\ref{sec:setup}), producing visibly underexposed video frames; all other positions use 3.91\,ms calibration with recovery ranging nonmonotonically from 33.8--69.4\% across geometry. The pooled value is therefore influenced by exposure calibration and does not represent the risk at any single position; deployment assessment should use worst-case geometry and exposure. SSIM from the synthetic-corpus digital pipeline reveals a proxy trade-off between security and reconstruction fidelity (Fig.~\ref{fig:tradeoff}); it is not real-capture or human-subject evidence. The protocol in \S\ref{sec:user_study} excludes the capture-hardened profile, leaving this profile's perceptual cost outside the current evidence.

\item \textbf{Descriptive trends across geometric conditions}: Across the tested 9~geometric conditions (0.5--1.5\,m $\times$ 0--30$^{\circ}$), protected-condition short-exposure recovery rates are consistently below corresponding unprotected values. Because the design is unbalanced and no inferential model for geometric factors has been established, this paper does not claim statistically significant monotonic robustness in either the distance or angle dimension.
\end{enumerate}

\begin{figure}[!t]
\centering
\includegraphics[width=\columnwidth]{figures/real_capture_montage.pdf}
\caption{Real-capture montage. Within each vertically stacked profile group (unprotected, deployed, and capture-hardened), panels show, from top to bottom, the digital integrated reconstruction view (for checking full-cycle image consistency), short-exposure capture, and long-exposure capture. Digital integrated views are not direct evidence of human readability or visual comfort.}
\label{fig:montage}
\end{figure}

\begin{figure}[!t]
\centering
\includegraphics[width=\columnwidth]{figures/real_capture_bar.pdf}
\caption{Best-of-engine OCR character recovery and exact match for three representative real-capture profiles---unprotected, deployed, and capture-hardened---under short exposure, long exposure, and video temporal averaging, pooled across the 9 tested distance--angle configurations. Lower bars indicate less recovered text. The long-exposure and video bars expose the integration boundary: the readability-priority deployed profile retains substantial recoverable content even when its short-exposure bar is low. Deployed sample counts include the disclosed five-item repeat and therefore differ from the other profiles.}
\label{fig:bar_attack}
\end{figure}

\begin{figure}[!t]
\centering
\includegraphics[width=\columnwidth]{figures/readability_robustness_tradeoff.pdf}
\caption{Trade-off between digital reconstruction quality and integration-attack robustness: as anti-OCR enhancement increases (no overlay $\to$ strong overlay $\to$ deployed (overlay + $\alpha=0.2$ inversion) $\to$ capture-hardened), temporal-average character recovery decreases (left axis) and full-cycle digital reconstruction SSIM~\cite{b_wang2004} also decreases (right axis). Data from anti-OCR profile ablation on synthetic corpus (digital pipeline, not real capture); SSIM is an image proxy, not a user readability measure.}
\label{fig:tradeoff}
\end{figure}

\textbf{Per-engine results}: Table~\ref{tab:real_ocr_engine} breaks down short-exposure results by engine; Fig.~\ref{fig:real_engine_ocr} visualizes all six profiles therein. The three engines differ substantially in unprotected character recovery (Surya 37.1\%, EasyOCR 79.5\%, Tesseract 84.3\%). Under the Strong profile, recovery ranges from 2.1--16.3\%; under deployed, 3.2--11.7\%; and under capture-hardened, 1.8--2.0\%. This shows recovery reduction across all three conventional OCR engines tested, but with limited sample sizes and engine families, the finding cannot be claimed as recognizer-agnostic. The main table's best-of-engine takes the highest recovery among three engines per capture, modeling an attacker free to choose the best recognizer. Its values are therefore no lower than any single-engine mean---a more conservative (i.e., less favorable to the defense) reporting caliber from the defender's perspective, not a better-performing one.

\begin{table}[!t]
\caption{Real-Capture Short-Exposure OCR: Per-Engine Results (Pooled Across 9 Geometric Conditions)}
\label{tab:real_ocr_engine}
\centering
\small
\setlength{\tabcolsep}{5.5pt}
\begin{tabular}{llccc}
\toprule
Profile & Engine & Char (\%) & Exact (\%) & N \\
\midrule
\multirow{3}{*}{Original}
 & Surya & 37.1 & 22.8 & 324 \\
 & EasyOCR & 79.5 & 27.8 & 324 \\
 & Tesseract & 84.3 & 47.8 & 324 \\
\midrule
\multirow{3}{*}{Mask only}
 & Surya & 4.2 & 0.3 & 324 \\
 & EasyOCR & 15.6 & 0.0 & 324 \\
 & Tesseract & 3.2 & 0.0 & 324 \\
\midrule
\multirow{3}{*}{Mask + noise}
 & Surya & 11.0 & 2.8 & 324 \\
 & EasyOCR & 18.1 & 0.0 & 324 \\
 & Tesseract & 4.8 & 0.0 & 324 \\
\midrule
\multirow{3}{*}{Strong}
 & Surya & 6.1 & 0.6 & 324 \\
 & EasyOCR & 16.3 & 0.0 & 324 \\
 & Tesseract & 2.1 & 0.0 & 324 \\
\midrule
\multirow{3}{*}{Deployed}
 & Surya & 3.2 & 0.4 & 459 \\
 & EasyOCR & 11.7 & 0.0 & 459 \\
 & Tesseract & 3.3 & 0.0 & 459 \\
\midrule
\multirow{3}{*}{Capture-hardened}
 & Surya & 1.8 & 0.0 & 324 \\
 & EasyOCR & 2.0 & 0.0 & 324 \\
 & Tesseract & 1.9 & 0.0 & 324 \\
\bottomrule
\end{tabular}
\end{table}

\begin{figure}[!t]
\centering
\includegraphics[width=\columnwidth]{figures/real_engine_ocr.pdf}
\caption{Real-capture short-exposure character recovery for three OCR engines across all six profiles in Table~\ref{tab:real_ocr_engine}; error bars are 95\% bootstrap intervals resampled by capture. All three engines are suppressed, but with a limited engine family this cannot be claimed as recognizer-agnostic.}
\label{fig:real_engine_ocr}
\end{figure}

\subsection{Real-Capture VLM Probes}
\label{sec:vlm}

To test whether conventional OCR conclusions can be extrapolated to stronger end-to-end vision recognizers, we applied three commercial vision-language models (Qwen3-VL-32B-Instruct, Kimi-K2.6, and GLM-4.5V, via SiliconFlow API) to real-capture photographs. This section is not designed to extend the protection claim; it is intentionally constructed as a \emph{boundary/counterexample probe}: if VLMs can recover text from images where conventional OCR fails, then the application value of this work must be explicitly bounded to conventional OCR short-exposure mitigation and cannot claim effectiveness against modern vision-language models.

The probes cover the same UVC camera at the 0$^{\circ}$ on-axis position across three distances. The close-range 0.5\,m run includes 156 input images/aggregate views: 36 same-session unprotected short-exposure baselines, 36 capture-hardened short/long-exposure each, and 4 types of aggregate views synthesized offline from 60\,fps video (12 segments each). Additional 1.0\,m and 1.5\,m runs each include 120 input images/aggregate views for the capture-hardened profile: 36 short/long-exposure each and 4~video aggregate view types (12 segments each). The VLM inputs across all three distances and the OCR BoE reference column in Table~\ref{tab:real_vlm} use the same batch of files (short-exposure/video-frame exposure 3.91\,ms, long exposure 31.25\,ms); reference values are computed on those files using the best-of-engine caliber per \S\ref{sec:setup}. All inputs are raw camera JPEGs without contrast enhancement, aligned with the OCR pipeline; pilot tests found that CLAHE-type local enhancement partially reconstructs masked text, effectively turning preprocessing into an additional attack component. The three runs correspond to 1,188 model requests, of which 1,133 returned parseable results: the 0.5\,m run had 34 failures out of 468 calls (Qwen~2, Kimi~0, GLM~32); the 1.0\,m run had 21 failures out of 360 (Qwen~4, Kimi~1, GLM~16); the 1.5\,m run had zero failures out of 360. Each cell reports its effective call count $N$. Accordingly, this paper uses the failure-free 1.5\,m run as the primary complete evidence for the additional distances, while the 1.0\,m run with failures is included in Table~\ref{tab:real_vlm}'s short-exposure distance sweep as descriptive supplementary data only. At 1.0\,m, Qwen and Kimi long-exposure exact are 84.4\% and 74.3\% respectively, and all three models' video temporal-average exact ranges from 70.0--83.3\%---high-recovery trends consistent with 1.5\,m. However, GLM long-exposure exact at 1.0\,m is only 3.3\%, markedly lower than 58.3\% at 1.5\,m; thus no cross-distance consistency generalization is made across all three models, and per-record results are in the data archive. For compactness, Table~\ref{tab:real_vlm} reports only the temporal-mean video aggregate at 0.5\,m as the representative full-cycle integration view; the other 0.5\,m aggregate views (single-best frame, best sliding-window mean, and max projection) are preserved in the data archive. The failure-free 1.5\,m run reports all four aggregate views to show view-selection sensitivity. Model identifiers are Qwen/Qwen3-VL-32B-Instruct, Pro/moonshotai/Kimi-K2.6, and zai-org/GLM-4.5V; decoding temperature is 0, maximum output 4,096 tokens. All three models share a single transcription prompt: transcribe only visible text, do not guess missing characters; reference ground truth is not included in the prompt. The full prompt, calling parameters, and per-call model outputs (transcriptions, self-reported confidence, and notes) are archived with the code and data (see ``Data and Code Availability''). Metrics are computed identically to \S\ref{sec:setup}, pooled at the capture level. CET6 documents are publicly available and may appear in model training data, potentially inflating character recovery through memorization; synthetic private snippets and exact-match are therefore the cleaner discriminative caliber.

\textbf{0.5\,m session selection transparency}: The archive also preserves a \path{real_capture_vlm_d0.5_a0_rerun.json} (file timestamp 012715) alternative session at the same position; its capture-hardened short-exposure inputs are nearly all-black, all three models' transcriptions are empty, and char/exact are both 0\%. This session is archived as an underexposure diagnostic but does not enter Table~\ref{tab:real_vlm} or get averaged with the main run. The paper uses the 141107, 3.91\,ms session because it shares the same file batch with the OCR reference column and has more adequate exposure---more advantageous to the attacker. The selection principle is input alignment and conservative reporting, not cherry-picking by recovery rate. Without independent photometer data, this paper treats ``underexposure'' only as a session-level diagnostic supported jointly by near-black images and empty transcriptions, not as a confirmed camera ISP causal conclusion.

\begin{table*}[!t]
\caption{Real-Capture VLM Probes (0$^{\circ}$ On-Axis; All Rows Except the Original Short Baseline Are Capture-Hardened). Each VLM Cell Reports Exact / Char / $N$ (\% / \% / Effective Calls); OCR BoE Column Shows Same-Batch Three-Engine Best-of-Engine Reference (Exact / Char, \%). For compactness, 0.5\,m reports only the temporal-mean video aggregate; the remaining 0.5\,m aggregate views are archived, while 1.5\,m reports all four aggregate views.}
\label{tab:real_vlm}
\centering
\footnotesize
\setlength{\tabcolsep}{4.2pt}
\begin{tabular}{llcccc}
\toprule
Distance & Condition & OCR BoE & Qwen3-VL-32B & Kimi-K2.6 & GLM-4.5V \\
\midrule
\multicolumn{6}{l}{\textit{Short-exposure distance sweep}}\\
\midrule
0.5\,m & original short & 58.3~/~92.4 & 83.3~/~94.5~/~36 & 83.3~/~95.6~/~36 & 73.1~/~92.0~/~26 \\
0.5\,m & short & 0.0~/~8.4 & 77.8~/~91.5~/~36 & 47.2~/~84.7~/~36 & 51.6~/~70.8~/~31 \\
1.0\,m & short & 0.0~/~3.6 & 55.6~/~79.1~/~36 & 38.9~/~68.6~/~36 & 27.3~/~47.0~/~33 \\
1.5\,m & short & 0.0~/~6.5 & 47.2~/~68.9~/~36 & 33.3~/~61.5~/~36 & 36.1~/~70.0~/~36 \\
\midrule
\multicolumn{6}{l}{\textit{Key attack views}}\\
\midrule
0.5\,m & long & 0.0~/~2.1 & 82.4~/~85.7~/~34 & 2.8~/~7.9~/~36 & 0.0~/~4.9~/~30 \\
0.5\,m & video:temporal\_mean & 8.3~/~50.7 & 83.3~/~94.0~/~12 & 83.3~/~95.5~/~12 & 87.5~/~91.8~/~8 \\
1.5\,m & long & 0.0~/~13.2 & 91.7~/~89.3~/~36 & 88.9~/~89.9~/~36 & 58.3~/~62.4~/~36 \\
1.5\,m & video:single\_best & 0.0~/~4.3 & 8.3~/~24.8~/~12 & 0.0~/~12.7~/~12 & 0.0~/~29.8~/~12 \\
1.5\,m & video:temporal\_mean & 0.0~/~48.4 & 66.7~/~90.3~/~12 & 66.7~/~90.3~/~12 & 58.3~/~88.5~/~12 \\
1.5\,m & video:window\_mean\_best & 0.0~/~8.6 & 58.3~/~89.0~/~12 & 41.7~/~84.2~/~12 & 66.7~/~87.2~/~12 \\
1.5\,m & video:max\_proj & 0.0~/~19.0 & 66.7~/~90.1~/~12 & 58.3~/~89.2~/~12 & 58.3~/~88.7~/~12 \\
\bottomrule
\end{tabular}
\end{table*}

\begin{table}[!t]
\caption{0.5\,m Capture-Hardened Short-Exposure VLM Character Recovery by Content Type. Baseline Column Shows Same-Session Unprotected Original Short Results for Qwen (Kimi/GLM Similar); GLM Had Occasional Call Failures, Parenthesized $N$ Shows Effective Count. At 1.5\,m, CET6 Dense Pages Yield 0\% Char Across All Three Models.}
\label{tab:vlm_content}
\centering
\footnotesize
\setlength{\tabcolsep}{2.8pt}
\begin{tabular}{lcccc}
\toprule
Content type ($N$) & Baseline (\%) & Qwen (\%) & Kimi (\%) & GLM (\%) \\
\midrule
CET6 document (3) & 64.4 & 18.5 & 5.6 & 0.4 \\
Credentials (6) & 97.6 & 97.6 & 88.3 & 77.1~(5) \\
Code snippets (6) & 89.8 & 94.9 & 95.8 & 88.9 \\
Digit strings (6) & 100.0 & 100.0 & 98.6 & 50.0~(4) \\
Mixed content (6) & 100.0 & 100.0 & 92.5 & 98.2~(4) \\
Paragraphs (6) & 100.0 & 99.8 & 81.7 & 66.5 \\
English sentences (3) & 95.3 & 95.3 & 96.9 & 94.6 \\
\bottomrule
\end{tabular}
\end{table}

Table~\ref{tab:real_vlm} presents the pooled VLM probe results: short-exposure rows show distance variation, while long-exposure and video aggregate rows show integration attack boundaries. Table~\ref{tab:vlm_content} breaks down 0.5\,m capture-hardened short-exposure character recovery by content type.

\textbf{Key findings}:

\begin{enumerate}
\item \textbf{VLM probes bound the claim rather than extending it}: On the same batch of photographs at 0.5\,m, three-engine best-of-engine on the hardened short-exposure condition yields only 8.4\% char / 0\% exact, while Qwen3-VL on the same images achieves 91.5\% char / 77.8\% exact. At 1.5\,m, the three models retain 33.3--47.2\% short-exposure exact / 61.5--70.0\% char. Even with a prompt prohibiting guessing, VLMs recover substantial text from visible fragments in synthetic credentials, digit strings, and short sentences. The conclusion of this section is therefore not that ``the capture-hardened profile works against stronger attackers'' but rather that conventional OCR claims cannot be extrapolated to VLM attackers.

\item \textbf{Failure concentrates on large-font short snippets; dense small text remains protected}: At 0.5\,m close range, CET6 full-page dense documents yield 0.4--18.5\% VLM character recovery under short exposure (unprotected baseline 64.4\%), whereas credentials, code, digit strings, and other large-font categories yield 50--100\% (Table~\ref{tab:vlm_content}). At 1.5\,m, CET6 dense pages yield 0\% char / 0\% exact across all three models, while digit strings and English short sentences maintain high recovery. This pattern is consistent with mask granularity: small-font strokes are approximately 1--2~pixels wide; removing $1-1/n$ of pixels destroys glyph topology. Large-font glyphs retain global contours even through sparse fragments and stripes, recoverable by robust visual encoding and language priors. Conventional OCR depends on binarization and character segmentation, which stripe distortion disrupts; VLMs bypass that bottleneck.

\item \textbf{Attack capability varies markedly with model and imaging conditions}: Under 0.5\,m long exposure, Qwen3-VL still recovers 82.4\% exact while Kimi/GLM do not exceed 2.8\%; at 1.5\,m long-exposure exact rises to 58.3--91.7\%. This variation suggests that VLM screen-reading ability is simultaneously influenced by model version, distance, exposure, and image content; any conclusion about ``defeating current VLMs'' will expire as models and service versions change.

\item \textbf{Video temporal averaging remains the VLM residual frontier}: At 0.5\,m temporal-average, three models achieve 83.3--87.5\% exact / $\geq$91.8\% char; at 1.5\,m temporal-average, 58.3--66.7\% exact / 88.5--90.3\% char. By contrast, 1.5\,m single-best-frame video view yields only 0--8.3\% exact, and conventional OCR on the same temporal-average views stays at or below 8.3\% exact in the two temporal-average rows tabulated here (48.4--50.7\% char). This shows that for VLM attackers, the full-cycle integration boundary (\S\ref{sec:security_frontier}) is more dangerous than selecting a single clear frame, and does not necessarily require precise alignment or complex de-striping post-processing.
\end{enumerate}

In summary, the effectiveness claim for the capture-hardened profile must be bounded by content and attacker type: it holds for conventional OCR pipelines and dense small-text content; for attackers equipped with modern VLMs, protection of large-font short snippets (credentials, code, digit strings) under short exposure has already failed. This negative result is part of this paper's boundary analysis: it repositions the practical application from ``blocking arbitrary machine screen-reading'' to ``reducing conventional OCR short-exposure batch recovery,'' and defines VLM-aware display protection as a subsequent research problem. VLMs and video temporal averaging jointly constitute the residual attack frontier; countermeasure directions are discussed in \S\ref{sec:vlm_defense}. This section covers only a single camera, the 0$^{\circ}$ on-axis position across three distances, a single protection profile, and three commercial API models, with only 12 video segments per cell and commercial API models that may update over time. Numbers should therefore be interpreted as descriptive boundary probe results and cannot substitute for full 9-geometry, cross-camera, or cross-model-version VLM evaluation.

% ----------------------------------------------------------
\subsection{User Experience Study}
\label{sec:user_study}

To evaluate the impact of the protected display on human readability and visual comfort, we designed a web-based user study.

\textbf{Equipment and environment protocol}: The same 240\,Hz display as in \S\ref{sec:setup}. The formal study version sets a hard measured refresh-rate threshold of $\geq$200\,Hz to ensure the fixed $n=4$ basic subframe cycle rate $f_r/n\geq50$\,Hz; below this threshold participation is blocked rather than silently reducing $n$. The 144--199\,Hz adaptive path is available only for an explicitly labeled demonstration mode, and its sessions are excluded from analysis by default. Before each participant begins, the experimenter confirms fixed display brightness, disabled auto-brightness/power-saving and variable refresh rate, unified browser fullscreen, closed high-load background processes, and an approximate 60\,cm viewing distance. Refresh rate is re-measured after any display-mode change.

The browser records actual frame intervals in each trial's \texttt{requestAnimationFrame} loop; if the interval between adjacent callbacks exceeds 1.5$\times$ the expected interval computed from the pre-trial measured refresh rate, a dropped-frame event is logged, along with estimated drop count, drop rate, mean/max frame intervals, observed refresh rate, basic subframe cycle rate, and full cycle rate with inversion frame. This recording reflects only browser presentation cadence and cannot substitute for panel scanning or photometric measurement.

\textbf{Ethics and safety}: The planned study involves minimal-risk human-subjects research. The authors' institution does not maintain an ethics review board for non-medical behavioral experiments of this type; accordingly, the protocol specifies informed consent and the safety measures described below, but no formal data collection has yet been conducted. Before starting, participants must separately check informed consent and photosensitivity screening: the page explicitly states that the masked display involves rapid temporal modulation, that participants should stop immediately upon experiencing discomfort, eye fatigue, dizziness, nausea, or headache, and that they may withdraw unconditionally at any time. Individuals who self-report a history of photosensitive epilepsy or sensitivity to flicker are not permitted to proceed. The system records consent timestamps and screening pass flags. Implementation-level safety is ensured by a $f_r/n\geq50$\,Hz floor; at $n=4$ on a 240\,Hz panel the basic cycle rate is 60\,Hz. The inversion frame reduces the full cycle frequency to $f_r/(n+1)=48$\,Hz; the implementation separately logs warnings and actual timing. Registration information (student ID, name) is used only for participation management; analysis and publication use only de-identified data. Age and gender are optional demographic fields.

\textbf{Task design}: Each participant first completes one 12-second, unscored unmasked Canvas typing warm-up, then views a 10-second, unscored deployed full-mask preview so that the first formal masked trial does not simultaneously serve to familiarize the stimulus, followed by four 20-second scored trials (within-subjects design). Two conditions are each repeated twice, counterbalanced using ABBA or BAAB order based on the experimenter's continuously assigned zero-based sequence number $k$ to control for practice and fatigue effects: (A)~\textbf{original text} (control), displayed as unmasked Canvas with $n=1$; (B)~\textbf{masked text}, using the actual deployment configuration: $n=4$, mask+noise, strong anti-OCR ($\alpha_s=0.10$, $\alpha_g=0.12$, 6 cycles) + weak inversion ($\alpha=0.2$). Specifically, the typing order index is $k\bmod2$ and the subjective rating Latin-square row is $\lfloor k/2\rfloor\bmod6$; every consecutive 12~participants thus fully cover the $2\times6$ joint assignment, ensuring deterministic and balanced joint order allocation. Formal sessions enforce a uniqueness constraint on $k$ in the database, and the backend recomputes and verifies both order indices. Both conditions are rendered by the same \texttt{renderSourceCanvas} path producing light text on a dark background, with identical canvas size, 23\,px bold font, and color---the sole processing difference is the temporal protection pipeline, thereby avoiding polarity and font-weight confounds. Here ``cycles'' refers to the web implementation's rotation of mask, noise, and stripe patterns through distinct instances before repeating; a value of 6 prevents long-exposure photography from integrating a single fixed pattern into a recognizable residue, sharing $\alpha_s$ and $\alpha_g$ parameters with the real-capture configuration described in \S\ref{sec:method} but constituting a web-playback-specific setting.

Scoring no longer compares characters position-by-position but aligns the full transcription against the optimal target prefix using Levenshtein minimum string distance (MSD)~\cite{b_soukoreff2003}; untyped target suffix does not count as error. From the alignment, accuracy, CPM, and WPM are computed, along with edit distance, aligned prefix length, and scoring version. First-keystroke latency from clicking ``Start'' to the first valid input event is also recorded. Analysis first averages each participant's two trials per condition, then performs within-subject paired comparison.

\textbf{Subjective ratings}: After the typing trials, participants rate 6~conditions on 4~dimensions using 1--5 Likert scales: readability, flicker perception, immediate visual discomfort (fatigue field), and perceived privacy. Conditions include the unmasked original anchor, $n\in\{2,3,4\}$ mask+noise, $n=4$ mask-only, and the same $n=4$ deployed full configuration (mask+noise+strong anti-OCR+weak inversion) used in the typing task. The first three masking levels all achieve $\geq$60\,Hz real temporal playback on a 240\,Hz panel; the originally planned $n=8$ would yield 30\,Hz and trigger static fallback, so it is not included as a temporal ablation. The six conditions are presented in a 6th-order balanced Latin square; each condition must be viewed for at least 10~seconds before submission is enabled, with viewing start/end times and duration recorded. The unmasked anchor also serves as an attention/scale-interpretation aid but is not mechanically excluded by a ``must rate 5'' rule. This study does not include the capture-hardened profile; assessment of its visible stripes and grain's acceptability is left to future work (see \S\ref{sec:discussion}).

\textbf{Sample size and analysis plan}: A minimum of $N=24$ participants is planned, a number within the range required for paired effect sizes of $d_z\approx0.6$--$0.8$ at 80\% power, and exactly covering two rounds of the 12-person joint order assignment. Pre-registered exclusion criteria: debug/demo sessions, incomplete four typing or six rating records, measured refresh rate below 200\,Hz, any typing trial with fewer than 5 attempted characters, control mean accuracy below 50\%, masked trials not in temporal mode, or any trial with observed basic cycle rate below 50\,Hz. WPM, CPM, accuracy, attempted characters, and first-keystroke latency are first averaged within each participant by condition; speed metrics must be interpreted jointly with accuracy to avoid misinterpreting fast input under low accuracy as performance improvement. Paired differences are tested using paired $t$-tests when the pre-specified Shapiro test is passed, otherwise Wilcoxon signed-rank tests, with effect sizes and participant-level bootstrap 95\% confidence intervals reported. Likert data are analyzed per dimension using Friedman tests with Holm-corrected post-hoc pairwise Wilcoxon tests. Analysis scripts generate de-identified participant means, JSON audit reports, and two \LaTeX\ tables directly from the formal database \texttt{study\_formal.db}.

\textbf{Participants}: \TODO{Final participant count and demographics.}

\TODO{Table: Typing trial objective metrics (accuracy/CPM/WPM, control vs.\ masked, paired differences + CI).}

\TODO{Table: 6 conditions $\times$ 4 dimensions Likert means + SD.}

\TODO{Core finding: whether the deployed full configuration produces significant changes in typing accuracy/speed and first-keystroke latency; subjective readability, flicker perception, and immediate visual discomfort rating results.}

% ----------------------------------------------------------
\subsection{Real-Capture Object Detection and Tracking}
\label{sec:detection}

To examine whether visual tasks beyond text recognition are affected by similar temporal degradation, this section reports an exploratory detection/tracking stress test within the same camera link, with evidence strength not equivalent to the OCR main experiment. Content was displayed on the 240\,Hz screen, captured by the eMeet S600 at 1.5\,m / $0^{\circ}$, and processed offline by four detectors~\cite{b_yolo26,b_rtdetr2024,b_fasterrcnn2015,b_retinanet2017}. Detection uses real\_clean (unprotected content captured through the same screen-camera link) as reference; results appear in Table~\ref{tab:real_coco}. Tracking uses 450~frames from MOT17-09-FRCNN with a per-source-frame still-display-then-capture verification procedure rather than continuous video synchronization; results appear in Table~\ref{tab:real_mot}. The associator is a greedy-fallback implementation inspired by ByteTrack~\cite{b_bytetrack2022}, not official ByteTrack. The real\_clean reference partially controls for screen-camera link degradation, but exposure interactions between the protected display and camera may still exist, so differences cannot be solely attributed to a single cause or extrapolated to other camera links.

\begin{table}[!t]
\caption{Real-Capture COCO Detection (150 Images; real\_clean = Unprotected Capture Baseline)}
\label{tab:real_coco}
\centering
\small
\begin{tabular}{llccc}
\toprule
Detector & Condition & mAP & mAP50 & AR \\
\midrule
\multirow{3}{*}{YOLO26x}
 & real\_clean & 28.6\% & 43.9\% & 31.0\% \\
 & real\_short & 1.3\% & 2.4\% & 1.4\% \\
 & real\_video & 3.5\% & 5.3\% & 3.6\% \\
\midrule
\multirow{3}{*}{RT-DETR-x}
 & real\_clean & 30.1\% & 50.2\% & 34.8\% \\
 & real\_short & 3.5\% & 5.8\% & 4.1\% \\
 & real\_video & 4.6\% & 7.3\% & 5.5\% \\
\midrule
\multirow{3}{*}{Faster R-CNN}
 & real\_clean & 21.3\% & 39.2\% & 28.7\% \\
 & real\_short & 0.8\% & 1.6\% & 0.9\% \\
 & real\_video & 1.5\% & 3.0\% & 1.8\% \\
\midrule
\multirow{3}{*}{RetinaNet}
 & real\_clean & 22.6\% & 39.4\% & 27.0\% \\
 & real\_short & 0.9\% & 1.8\% & 1.0\% \\
 & real\_video & 1.9\% & 3.1\% & 2.0\% \\
\bottomrule
\end{tabular}
\end{table}

\begin{table}[!t]
\caption{Real-Capture MOT17-09-FRCNN Still-Frame Tracking (ByteTrack-Style Greedy Association, 450 Frames)}
\label{tab:real_mot}
\centering
\begin{tabular}{llcc}
\toprule
Detector & Condition & HOTA & IDF1 \\
\midrule
\multirow{3}{*}{YOLO26x}
 & real\_clean & 15.8\% & 12.1\% \\
 & real\_short & 8.7\% & 6.5\% \\
 & real\_video & 10.0\% & 8.7\% \\
\midrule
\multirow{3}{*}{RT-DETR-x}
 & real\_clean & 14.2\% & 10.4\% \\
 & real\_short & 8.1\% & 5.2\% \\
 & real\_video & 9.6\% & 6.8\% \\
\midrule
\multirow{3}{*}{Faster R-CNN}
 & real\_clean & 14.4\% & 9.6\% \\
 & real\_short & 6.8\% & 6.0\% \\
 & real\_video & 10.5\% & 7.1\% \\
\midrule
\multirow{3}{*}{RetinaNet}
 & real\_clean & 14.2\% & 9.8\% \\
 & real\_short & 5.1\% & 4.4\% \\
 & real\_video & 6.8\% & 5.9\% \\
\bottomrule
\end{tabular}
\end{table}

Under short exposure, all four detectors' mAP50 drops from the unprotected real-capture baseline of 39.2--50.2\% to 1.6--5.8\% (Fig.~\ref{fig:real_det_drop}), representing relative decreases exceeding 85\%. The real\_video condition yields 3.0--7.3\%, also below respective real\_clean baselines. This pattern is consistent with visual feature degradation under the protected condition and shows that the single-frame destruction is not limited to character-segmentation-based OCR. However, with only one camera, one detection position, and 150~images, this paper does not extrapolate the finding to a cross-device or general object-detection protection conclusion. In still-frame tracking, the four detectors' HOTA drops from 14.2--15.8\% to 5.1--8.7\% under short exposure, and IDF1 drops from 9.6--12.1\% to 4.4--6.5\%. It should also be noted that the real\_clean baseline HOTA (14.2--15.8\%) is itself substantially depressed by the still-frame screen-capture link, far below typical levels for this sequence on native images, compressing the remaining room for the protected condition to decrease further. These tracking results therefore provide only directional evidence. The still-frame procedure also cannot substitute for a continuous-video tracking experiment.

\begin{figure}[!t]
\centering
\includegraphics[width=\columnwidth]{figures/real_detection_drop.pdf}
\caption{Real-capture detection: mAP50 drop for four detectors under protected conditions.}
\label{fig:real_det_drop}
\end{figure}

% ----------------------------------------------------------
\subsection{Software Simulation Supplementary Experiments}
\label{sec:simulation}

The preceding experiments provide results under the current physical capture setup. This section uses software simulation to address two additional questions: (1)~If the attacker can obtain only a single protected subframe from the display compositor, what is its machine readability? (2)~What is the attacker's recovery boundary under zero camera noise with perfect alignment? Question~(1) also characterizes the conventional single-frame screenshot scenario: a screenshot tool captures the framebuffer at a single instant, corresponding to one protected subframe. The 0--2.6\% single-subframe recovery rates below therefore also constitute mitigation evidence against single-frame screenshots. This mitigation has a clear boundary: if malware intercepts the original frame before masking, or if screen recording captures and accumulates frames spanning a full cycle, the method cannot provide protection. Accordingly, this paper claims only single-frame-level mitigation for the screenshot scenario, not a general ``anti-screenshot'' guarantee.

\subsubsection{Synthetic Single-Subframe OCR}
Using Tesseract~\cite{b_smith2007}, EasyOCR~\cite{b_easyocr2020}, and Surya~\cite{b_surya2024} on 120 synthetic samples, character recovery drops from 94.0--97.0\% to 0--2.6\% across the three engines (Table~\ref{tab:ocr_corpus}). This result covers only the current corpus, parameters, and three engines. Stratified results can be used to check inter-category trends but are insufficient to prove effectiveness across all languages, layouts, or recognizers.

\begin{table}[!t]
\caption{Synthetic Single-Subframe OCR: Three-Engine Character Recovery (120 Samples, $n=4$, $\varepsilon=8/255$)}
\label{tab:ocr_corpus}
\centering
\begin{tabular}{lccc}
\toprule
OCR engine & Original frame & Single subframe & Reduction \\
\midrule
Tesseract & 94.0\% & 0.0\% & 94.0\% \\
EasyOCR & 94.1\% & 0.0\% & 94.1\% \\
Surya & 97.0\% & 2.6\% & 94.4\% \\
\bottomrule
\end{tabular}
\end{table}

\begin{figure}[!t]
\centering
\includegraphics[width=\columnwidth]{figures/multiengine_ocr.pdf}
\caption{Character recovery for three OCR engines on original frames and single protected subframes; error bars are 95\% bootstrap intervals resampled by corpus sample.}
\label{fig:multiengine_ocr}
\end{figure}

Fig.~\ref{fig:multiengine_ocr} visualizes the three-engine results from Table~\ref{tab:ocr_corpus}. All three engines' single-subframe character recovery does not exceed 2.6\% on the current data, indicating the result is not driven by a single engine. Broader recognizer generalization still requires additional experiments.

\subsubsection{Theoretical Security Boundary}
\label{sec:security_frontier}
When an attacker under ideal conditions accumulates $\geq n$ frames covering a full cycle, the complementarity is exactly exploited and adversarial noise is simultaneously cancelled per~(\ref{eq:noise_complement}), recovering 95.2\% (per-sample strongest attack). This is the \emph{fundamental boundary} of any linear temporal scheme. The capture-hardened profile reduces recovery under the current real-capture integration conditions, but video averaging still yields 47.9\% character recovery; it can therefore only be called mitigation, not closure of this boundary.

\subsubsection{Comparison with Software-Proxy Baselines}
In an exploratory simulation on 9~stratified samples, character recovery rates for dimming, Gaussian blur, pixelation, off-axis privacy-filter proxy, and a single temporal-mask subframe are 92.9\%, 52.6\%, 0\%, 84.7\%, and 0\%, respectively. Dimming by 50\% matches the unprotected condition at 92.9\%, indicating this parameter does not reduce current OCR recovery. These treatments' parameters, physical conditions, and user experience are not equivalently matched---in particular, the ``privacy filter'' is a software off-axis proxy, not a physical product test. Fig.~\ref{fig:baselines_scatter} therefore shows only the trade-off under the current proxy settings and does not support a conclusion of superiority over all passive approaches.

\begin{figure}[!t]
\centering
\includegraphics[width=\columnwidth]{figures/baselines_scatter.pdf}
\caption{Software-proxy baselines on 9~samples: horizontal axis is OCR character recovery, vertical axis is digital reconstruction $\Delta E_{00}$. Method parameters and physical conditions are not equivalently matched; the figure serves only as an exploratory comparison.}
\label{fig:baselines_scatter}
\end{figure}

\begin{figure*}[!t]
\centering
\includegraphics[width=0.82\textwidth]{figures/antiocr_heatmap.pdf}
\caption{Anti-OCR profile stripe $\times$ glyph amplitude grid (temporal-average attack): left panel shows character recovery, right panel shows digital reconstruction drift. High stripe and glyph amplitude regions descriptively reduce character recovery to $\sim$11\% across the current 12 parameter grid points (16 samples per point) while increasing reconstruction cost. No significance testing was performed.}
\label{fig:antiocr_heatmap}
\end{figure*}

The stripe/glyph parameter trade-off for the anti-OCR profile is shown in Fig.~\ref{fig:antiocr_heatmap}.

\subsubsection{Detection and Tracking Simulation}
Supplementary simulations were run on COCO val2017~\cite{b_coco2014} and MOT17~\cite{b_mot16}. The four detectors are YOLO26x~\cite{b_yolo26}, RT-DETR-x~\cite{b_rtdetr2024}, Faster R-CNN~\cite{b_fasterrcnn2015}, and RetinaNet~\cite{b_retinanet2017}. The COCO simulation used only 8~images, positioned as implementation pipeline diagnostics (Table~\ref{tab:coco_sim}). The MOT simulation covers 7~sequences totaling 5,316~frames, using a ByteTrack-inspired~\cite{b_bytetrack2022} greedy association fallback implementation (Table~\ref{tab:mot_sim}). Because TrackEval failed to run, MOTA/MOTP/IDF1 in the table come from an in-project approximate evaluation backend and are not equivalent to official ByteTrack/TrackEval results.

\begin{table}[!t]
\caption{COCO val2017 Simulated Detection Pipeline Diagnostics (4 Detectors, 8 Images; Not a Dataset-Level Performance Estimate)}
\label{tab:coco_sim}
\centering
\small
\setlength{\tabcolsep}{5.8pt}
\begin{tabular}{llccc}
\toprule
Detector & Condition & mAP & mAP50 & AR \\
\midrule
\multirow{3}{*}{YOLO26x}
 & clean & 48.1\% & 58.3\% & 51.0\% \\
 & single subframe & 0.0\% & 0.0\% & 0.0\% \\
 & temporal avg. & 14.3\% & 18.3\% & 14.3\% \\
\midrule
\multirow{3}{*}{RT-DETR-x}
 & clean & 47.1\% & 58.7\% & 53.4\% \\
 & single subframe & 5.2\% & 6.0\% & 5.2\% \\
 & temporal avg. & 16.8\% & 23.7\% & 17.9\% \\
\midrule
\multirow{3}{*}{Faster R-CNN}
 & clean & 38.2\% & 52.7\% & 43.6\% \\
 & single subframe & 3.0\% & 4.4\% & 3.0\% \\
 & temporal avg. & 11.1\% & 16.5\% & 11.1\% \\
\midrule
\multirow{3}{*}{RetinaNet}
 & clean & 33.2\% & 47.1\% & 35.3\% \\
 & single subframe & 3.3\% & 4.2\% & 3.3\% \\
 & temporal avg. & 9.6\% & 13.2\% & 9.6\% \\
\bottomrule
\end{tabular}
\end{table}

\begin{table}[!t]
\caption{MOT17 Simulated Tracking (ByteTrack-Style Greedy Association; In-Project Approximate Metrics; 5,316 Frames)}
\label{tab:mot_sim}
\centering
\begin{tabular}{llccc}
\toprule
Detector & Condition & MOTA & MOTP & IDF1 \\
\midrule
\multirow{3}{*}{YOLO26x}
 & clean & 31.1\% & 81.7\% & 27.1\% \\
 & single subframe & 0.1\% & 74.8\% & 0.1\% \\
 & temporal avg. & 9.1\% & 78.1\% & 12.0\% \\
\midrule
\multirow{3}{*}{RT-DETR-x}
 & clean & 24.1\% & 79.8\% & 38.2\% \\
 & single subframe & $-$2.3\% & 73.2\% & 5.4\% \\
 & temporal avg. & $-$11.4\% & 75.0\% & 14.7\% \\
\midrule
\multirow{3}{*}{Faster R-CNN}
 & clean & 3.8\% & 77.9\% & 35.6\% \\
 & single subframe & $-$2.5\% & 71.2\% & 4.5\% \\
 & temporal avg. & $-$2.7\% & 73.2\% & 14.0\% \\
\midrule
\multirow{3}{*}{RetinaNet}
 & clean & $-$5.4\% & 77.1\% & 30.5\% \\
 & single subframe & 0.4\% & 70.6\% & 2.7\% \\
 & temporal avg. & $-$2.9\% & 73.8\% & 11.9\% \\
\bottomrule
\end{tabular}
\end{table}

On these 8~images, single-subframe mAP50 across four detectors is 0--6.0\%. This sample size is sufficient only to check trends, not to estimate COCO population performance. In the 5,316-frame approximate tracking results, single-subframe IDF1 is 0.1--5.4\%, though MOTA exhibits negative values and nonmonotonic behavior. Fig.~\ref{fig:all_attackers} summarizes residual capability of a single protected subframe relative to the clean baseline across the current evaluation; it does not represent untested models.

\begin{figure}[!t]
\centering
\includegraphics[width=\columnwidth]{figures/all_attackers.pdf}
\caption{Residual capability of a single protected subframe relative to the clean baseline (clean = 100\%) across the current evaluation: OCR, detection, and tracking are approximately 0--9.7\%. This figure summarizes only the tested models and metrics.}
\label{fig:all_attackers}
\end{figure}

\subsubsection{Security--Usability Trade-off}
The model sweep covers $n\in\{2,4,6,8\}$ and refresh rates $\in\{144,240,360\}$\,Hz across 12~combinations. The Pareto front is defined on the (FPI, normalized mutual information) bi-objective minimization; the selection rule within the candidate set is: first retain safety candidates with $\mathrm{FPI}<0.1$ (configurations exceeding this threshold are excluded even if their normalized mutual information is lower---e.g., $n=8$ @ 360\,Hz with FPI\,=\,0.219 is eliminated for exceeding the threshold), then select the candidate with the lowest normalized mutual information. This rule selects $n=6$ @ 360\,Hz (FPI\,=\,0.033, normalized mutual information proxy = 0.249, $\Delta E_{00}=0.60$, SSIM\,=\,0.9995). For comparison, the base configuration $n=4$ @ 240\,Hz has FPI\,=\,0.030, normalized mutual information proxy = 0.377, and $\Delta E_{00}=0.34$; the $n=6$ candidate trades higher FPI and $\Delta E_{00}$ (0.60) and lower SSIM for lower single-subframe information retention. This is only a candidate configuration selected by a simplified digital model; it has not been validated on a physical 360\,Hz display or in a user study and is therefore not claimed as a deployment recommendation. Fig.~\ref{fig:pareto} shows the Pareto front across different $n$ and refresh-rate combinations on the two proxy axes.

\begin{figure}[!t]
\centering
\includegraphics[width=\columnwidth]{figures/pareto_sweep.pdf}
\caption{Security--usability proxy Pareto front in the digital model: $n \times$ refresh rate. Neither FPI nor normalized mutual information is a direct measure of user experience or real attack success rate.}
\label{fig:pareto}
\end{figure}


%% ============================================================
%% VI. DISCUSSION
%% ============================================================
\section{Discussion}
\label{sec:discussion}

\subsection{Evidence Hierarchy and Claim Scope}

The most directly supported conclusion of this paper is: as a single-UVC-camera feasibility study, temporal pixel masking substantially reduces conventional OCR attackers' character recovery and exact-match rates from short-exposure screen photographs under the current eMeet S600 camera, display, corpus, and 9~geometric conditions. The 10,575 real-capture images reflect component, parameter, and attack-condition capture coverage within this UVC screen-camera link. Because the corpus is not a balanced full-factorial design and covers only one webcam, this paper does not interpret the total sample count as cross-device or cross-scenario statistical inference evidence. Detection/tracking results, VLM probes, and software simulations respectively answer ``does this degradation transfer to other visual tasks,'' ``can stronger recognizers break through,'' and ``where is the full-cycle integration boundary''---collectively defining boundaries rather than extending the primary validity claim.

\subsection{The Mathematical Tension of Linear Integration}

This scheme exploits the difference between human temporal integration and single-frame camera sampling. However, human temporal integration is fundamentally linear summation ($\sum\mathbf{I}_k = \mathbf{I}$), and a camera performing the same multi-frame accumulation can recover the original information. This is the fundamental boundary shared by all persistence-of-vision-based temporal schemes. Additionally, research on adversarial robustness of vision-language models~\cite{b_li2025} shows that large models can still extract partial semantic information from perturbed images; the real-capture probes in \S\ref{sec:vlm} confirm this risk: VLMs can recover large-font content from single short-exposure photographs and obtain even stronger signals from temporal-average and long-exposure views. Section~\ref{sec:security_frontier} honestly quantifies the integration boundary: full-cycle integration recovers 95.2\%.

\subsection{Role and Cost of the Capture-Hardened Profile}

The capture-hardened profile breaks strict completeness through stripe and glyph perturbations, reducing recovery under the current full-cycle integration conditions. Its digital reconstruction metrics simultaneously worsen, indicating this enhancement is not cost-free. In \S\ref{sec:real_capture}, this profile achieves 0\% short-exposure exact-match and 5.6\% video temporal-average exact-match. The latter's character recovery remains at 47.9\%, so it can only be characterized as mitigation rather than complete blocking. Whether stripes and grain are acceptable to users must be answered by user experiments; the study designed in \S\ref{sec:user_study} covers only the deployed profile and core mask conditions and does not include the capture-hardened profile. Its acceptability therefore remains an unanswered question within this paper's scope, and this paper does not prejudge it.

\subsection{VLM Attackers and Countermeasure Directions}
\label{sec:vlm_defense}

Section~\ref{sec:vlm} shows that protection failure concentrates on large-font short snippets, and that temporal averaging/long exposure further amplifies the signal available to VLMs: glyph contour fragments retained in a single subframe are sufficient for language-prior completion, and more-cycle coverage approaches the original text. This paper retains this result as one of its primary boundary contributions rather than relegating it to an appendix anomaly, because it directly determines the method's deployability narrative: the current scheme can serve as a mitigation mechanism for conventional OCR short-exposure capture but cannot serve as modern VLM screen-reading protection. Potential countermeasure directions include: (1)~\textbf{font-size-adaptive mask granularity}, scaling the mask cell with rendered stroke width so that a single subframe no longer preserves character topology; (2)~\textbf{complementary glyph decoys}, injecting false strokes and characters within the $\sum_{k}\mathbf{N}_k=\mathbf{0}$ framework that cancel within a cycle, causing single-frame attackers to obtain high-confidence \emph{incorrect} text rather than partially correct text---directly countering prior-based completion; (3)~\textbf{deployment-layer content policies}, rendering high-risk fields (credentials, etc.) in dense small font or enabling the hardened profile per field. None of these directions have been experimentally validated. Adversarial perturbations targeting VLM visual encoders are additionally constrained by digital-to-physical transferability (\S\ref{sec:related}), with expected limited benefit. More generally, information visible to the observer within the integration window is in principle recoverable by a camera covering a full cycle paired with a strong-prior model; VLMs reduce the fragment quantity needed for recovery. The defense objective should therefore shift from ``unreadable to any attacker'' toward increasing attack cost and preventing silent bulk collection.

\subsection{Deployment Path}

Real-world deployment would likely require OS driver-level framebuffer injection (DXGI / KMS-DRM / CoreDisplay), coordination with panel/backlight timing, and high-refresh panels meeting the target cycle rate. These steps introduce color management, luminance saturation, frame drops, and synchronization errors that could alter the algorithm's behavior in a real link---they cannot be treated as simple engineering packaging that leaves the conclusions unchanged. More importantly, the deployed profile's 96.5\% sensitive-token recovery under long exposure means that credentials, account identifiers, and comparable high-risk fields should not rely on that profile alone. Such fields should use the capture-hardened profile or dense small-font rendering under a field-level content policy, while retaining the explicit caveat that these options impose unresolved readability costs and do not close the video/VLM boundary.

\subsection{Ethics Statement}

This research is defensive in nature, aiming to protect screen content from unauthorized photography. The attack analysis (\S\ref{sec:security_frontier}) is intended for honest assessment of security boundaries, not as an attack guide. The planned user study would involve human subjects in minimal-risk research; its protocol specifies informed consent, photosensitivity screening, and voluntary withdrawal as described in \S\ref{sec:user_study}.

\subsection{Limitations}

\begin{enumerate}
\item The user study described in \S\ref{sec:user_study} is a protocol rather than completed evidence: no formal participant dataset is currently available, so the readability-priority deployed profile's usability remains unvalidated. When conducted, the study will use a controlled 240\,Hz laboratory environment; external validity will still require broader participant samples and display conditions.
\item Immediate visual discomfort and perceived privacy in the user study are both short-exposure single-item subjective scales. The former cannot be extrapolated to long-term visual fatigue; the latter primarily reflects a ``looks hard to photograph clearly'' impression and may be influenced by research-purpose priming, and cannot serve as evidence of objective anti-capture effectiveness or attack success rate.
\item The current PoC is application-layer and has not integrated OS drivers or hardware backlight. Luminance, HDR, and color-difference conclusions come from digital models, lacking photometer, panel-response, and high-speed camera measurements.
\item Real-capture OCR used only one eMeet S600 UVC webcam; ambient illuminance was not recorded; exposure values have two calibration sets. The dataset contains repeated content with an unbalanced design; bootstrap intervals cannot substitute for cross-device, cross-scenario replication. The physical evaluation should therefore be interpreted as single-UVC-camera feasibility evidence, not generalization to smartphones, smart glasses, or other camera ISPs.
\item Real-capture detection used only 150 COCO images at a single position, positioned as an exploratory cross-task stress test. Real-capture tracking is 450 still-frame captures from a single MOT17 sequence, not a continuous video synchronization experiment. Simulated detection used only 8~images; tracking uses in-project greedy association and approximate metrics.
\item VLM evaluation is a single-UVC-camera, 0$^{\circ}$-on-axis, single-profile, three-model adversarial probe; each cell reports effective call count $N$, which varies with model parseable returns. Video conditions have only 12 segments per cell, and commercial API models may update over time. The OCR reference column in Table~\ref{tab:real_vlm} uses the same file batch as VLM inputs; evaluation is still limited to 0$^{\circ}$ on-axis at three distances and cannot substitute for full 9-position VLM evaluation. Public document content may have inflated character recovery through training-data memorization. This paper did not evaluate multi-cycle adaptive reconstruction, rolling-shutter phase search, or malware reading the original frame before masking. Full-cycle integration remains the fundamental boundary.
\end{enumerate}


%% ============================================================
%% VII. CONCLUSION
%% ============================================================
\section{Conclusion}
\label{sec:conclusion}

This paper implemented and evaluated a temporal pixel masking method based on temporal integration differences within the tested screen-to-camera link. The primary objective is to reduce conventional OCR recovery under short-exposure capture while allowing human observers to integrate content across subframes. This objective is not equivalent to ``unreadable to any camera'' or ``effective against all machine vision models''; full-cycle exposure, video reconstruction, and modern VLMs can still recover partial information.

Within the unbalanced 10,575-image real-capture corpus composed of component, parameter, and attack experiments, the core short-exposure OCR subset's best-of-engine character recovery dropped from 94.1\% to 15.1\% (deployed profile) and 5.0\% (capture-hardened profile), with exact-match falling from 65.7\% to 0.4\% / 0\%. In controlled synthetic evaluation, three OCR engines' single-subframe character recovery was 0--2.6\%. Exploratory cross-task stress testing showed that on a single-position real-capture COCO subset from the same UVC camera link, four detectors' mAP50 dropped from 39.2--50.2\% to 1.6--5.8\%, though this result does not support a cross-device object-detection generalization claim. The base configuration's digital proxies are FPI\,=\,0.030 and $\Delta E_{00}=0.34$, which cannot substitute for user or photometric experiments. Video temporal averaging remains the residual attack frontier: the capture-hardened profile yields 47.9\% char / 5.6\% exact, not blocking this attack. The 0$^{\circ}$ on-axis VLM probes further show that modern vision-language models raise the hardened profile's large-font short-snippet short-exposure exact-match from conventional OCR's 0\% back to as high as 77.8\%; at 1.5\,m, short-exposure exact remains 33.3--47.2\%, video temporal-average exact is 58.3--66.7\%. This result is a boundary finding rather than protection success evidence: it indicates that VLM-aware display-side protection remains an unsolved problem and further constrains the applicability scope of the candidate mechanism presented here.

These results support temporal pixel masking as a candidate mitigation mechanism for conventional OCR short-exposure screen capture in the tested camera link. Its real-world usability, cross-device generalization, applicability to smartphone/smart-glasses camera links, robustness to full-cycle reconstruction, and robustness to VLM attackers all require further validation.

%% ============================================================
%% DATA AND CODE AVAILABILITY
%% ============================================================
\section*{Data and Code Availability}

The software implementation, experimental configurations, per-sample conventional OCR results, and VLM probe complete prompts, calling parameters, and per-call model outputs are planned for release as a de-identified research archive accompanying the paper. For archive traceability, the legacy \texttt{vlm} profile tag in historical real-capture records denotes the capture-hardened configuration defined in \S\ref{sec:hardened}. If the planned user study is conducted, its formal database \texttt{study\_formal.db} will contain student IDs, names, and other participation-management information and will not be directly released; public materials will include only de-identified participant-level results, exclusion audits, and reproducible analysis scripts. \TODO{Final archive URL, version number, and persistent identifier to be added.}

%% ============================================================
%% ACKNOWLEDGMENT
%% ============================================================
\section*{Acknowledgment}
\TODO{Acknowledgment to be filled.}\EOD


%% ============================================================
%% REFERENCES
%% ============================================================
\clearpage
\bibliographystyle{IEEEtran}
\bibliography{refs}

\vfill

\end{document}
